\documentclass[11pt]{article}
\usepackage{amsmath,amssymb}
\usepackage{geometry}
\usepackage{enumitem}
\usepackage{multicol}
\geometry{margin=0.8in}
\setlength{\parindent}{0pt}
\setlength{\columnsep}{0.3in}

\begin{document}

\begin{center}
  {\Large\textbf{Calculus: Differentiation and Continuity}}\\[3pt]
  {\large Challenging Practice --- 45 Minutes}
\end{center}

Name: \rule{2.7in}{0.4pt}\hfill Date: \rule{1.5in}{0.4pt}

\section*{Instructions}
Complete all 12 problems without a calculator. Allow about 22 minutes for
Problems 1--6 and 23 minutes for Problems 7--12. Show derivative steps,
limit calculations, and theorem hypotheses where requested. Give exact answers;
factored derivatives are acceptable. All angles are in radians.
Do not use L'H\^opital's rule or Taylor series. This cover sheet may be used
throughout the practice.

\section*{Basic Formulas and Facts}
\begin{multicols}{2}
\small
\textbf{Derivative definition and rules}
\[
  f'(a)=\lim_{h\to0}\frac{f(a+h)-f(a)}{h}
\]
\[
\begin{aligned}
  (u+v)'&=u'+v', & (cu)'&=cu',\\
  (u^n)'&=nu^{n-1}u',\\
  (uv)'&=u'v+uv',\\
  \left(\frac{u}{v}\right)'&=\frac{u'v-uv'}{v^2},\quad v\ne0,\\
  [f(g(x))]'&=f'(g(x))g'(x).
\end{aligned}
\]
Here $u$ and $v$ are differentiable functions of $x$, and $c$ is constant.

\medskip
\textbf{Common derivatives}
\[
\begin{aligned}
  (e^u)'&=e^u u', & (\ln u)'&=\frac{u'}{u}\ (u>0),\\
  (\sin u)'&=\cos(u)u',\\
  (\cos u)'&=-\sin(u)u',\\
  (\tan u)'&=\sec^2(u)u'.
\end{aligned}
\]

\textbf{Implicit differentiation}
\[
  \frac{d}{dx}(y^n)=ny^{n-1}y',\qquad
  \frac{d}{dx}(xy)=y+xy'.
\]
For $F(x,y)=C$, if $F_y\ne0$,
\[
  y'=-\frac{F_x}{F_y}.
\]
Here $F_x$ and $F_y$ mean differentiation with respect to one variable while
holding the other fixed. A horizontal tangent requires $F_x=0$, $F_y\ne0$;
a vertical tangent requires $F_y=0$, $F_x\ne0$.

\medskip
\textbf{Tangent line and second derivative}
\[
  y-f(a)=f'(a)(x-a),\qquad y''=\frac{d}{dx}(y').
\]

\columnbreak
\textbf{Limits}

Use limit laws when the required limits exist; a quotient law requires a
nonzero denominator limit. For continuous functions, substitute directly.
\[
\begin{gathered}
  \lim_{t\to0}\frac{\sin t}{t}=1,\qquad
  \lim_{t\to0}\frac{\tan t}{t}=1,\\
  \lim_{t\to0}\frac{1-\cos t}{t^2}=\frac12.
\end{gathered}
\]
\[
  \lim_{x\to a}f(x)=L
  \iff \lim_{x\to a^-}f(x)=\lim_{x\to a^+}f(x)=L.
\]

\textbf{Continuity and differentiability}

Continuity at $a$ requires $f(a)$ to be defined and
\[
  \lim_{x\to a^-}f(x)=\lim_{x\to a^+}f(x)=f(a).
\]
For a piecewise join, differentiability requires continuity and equal,
finite one-sided difference-quotient limits. Differentiability implies
continuity; continuity does not imply differentiability. Existence of
$f'(a)$ does not require $f'$ to be continuous at $a$.

\medskip
\textbf{Intermediate Value Theorem (IVT)}

If $f$ is continuous on $[a,b]$ and $N$ lies strictly between $f(a)$ and
$f(b)$, then $f(c)=N$ for some $c\in(a,b)$. This does not guarantee uniqueness.

\medskip
\textbf{Mean Value Theorem (MVT)}

If $f$ is continuous on $[a,b]$ and differentiable on $(a,b)$, then some
$c\in(a,b)$ satisfies
\[
  f'(c)=\frac{f(b)-f(a)}{b-a}.
\]

\textbf{Rolle's theorem}

Under the MVT hypotheses, if $f(a)=f(b)$, then $f'(c)=0$ for some
$c\in(a,b)$.

\medskip
A failed hypothesis means the theorem cannot be invoked; its conclusion
might still happen to hold.
\end{multicols}

\newpage
\begin{center}
  {\large\textbf{Part I: Chain Rule and Implicit Differentiation}}\\[2pt]
  Problems 1--6: approximately 22 minutes total
\end{center}

\begin{enumerate}[leftmargin=*,itemsep=0pt]
  \item Differentiate.
  \[
    y=\frac{e^{\sin(3x)}\ln(1+x^2)}{\sqrt{2+x^4}}
  \]
  \vspace{\stretch{1}}

  \item Differentiable functions $f$ and $g$ satisfy
  \[
    g(0)=2,\qquad g'(0)=-1,\qquad f(5)=3,\qquad f'(5)=4.
  \]
  For
  \[
    H(x)=\ln\!\left(1+\left[f\!\left(g(x)^2+1\right)\right]^2\right),
  \]
  find $H'(0)$ and the tangent line at $x=0$.
  \vspace{\stretch{1}}

  \item For $y=\sin(x^2)$, find $y''$ and evaluate it at
  $x=\sqrt{\pi/2}$.
  \vspace*{\stretch{1}}
\end{enumerate}

\newpage
\begin{center}
  {\large\textbf{Part I --- Continued}}
\end{center}

\begin{enumerate}[leftmargin=*,itemsep=0pt,resume]
  \item For the curve
  \[
    e^{xy}+x^2-y^2=2,
  \]
  find $dy/dx$ in terms of $x$ and $y$, and find the tangent line at $(1,0)$.
  \vspace{\stretch{1}}

  \item For the curve
  \[
    x^2y+y^3=2,
  \]
  find $dy/dx$ in terms of $x$ and $y$, and find $d^2y/dx^2$ at $(1,1)$.
  \vspace{\stretch{1}}

  \item Find all points on
  \[
    x^2+xy+y^2=12
  \]
  where the tangent is horizontal or vertical. Justify each classification
  and check that the points lie on the curve.
  \vspace*{\stretch{1}}
\end{enumerate}

\newpage
\begin{center}
  {\large\textbf{Part II: Limits, Continuity, and Theorems}}\\[2pt]
  Problems 7--12: approximately 23 minutes total
\end{center}

\begin{enumerate}[leftmargin=*,itemsep=0pt,resume]
  \item Evaluate both limits, showing the transformations used.
  \[
    \text{(a)}\quad \lim_{x\to0}\frac{\sqrt{1+4x}-1}{\sin(3x)}
    \qquad
    \text{(b)}\quad \lim_{x\to0}\frac{1-\cos(4x)}{x\sin(2x)}
  \]
  \vspace{\stretch{1}}

  \item Determine all real values of $a$, $b$, and $k$ for which $p$ is
  differentiable at $x=2$. Then find $p'(2)$. Show the continuity and
  differentiability conditions separately.
  \[
    p(x)=
    \begin{cases}
      \dfrac{x^2-4}{x-2}, & x<2,\\[6pt]
      k, & x=2,\\[2pt]
      ax^2+b, & x>2.
    \end{cases}
  \]
  \vspace{\stretch{1}}

  \item Let
  \[
    q(x)=
    \begin{cases}
      x^2\sin(1/x), & x\ne0,\\
      0, & x=0.
    \end{cases}
  \]
  Determine whether $q$ is continuous at $0$, whether $q'(0)$ exists, and
  whether $q'$ is continuous at $0$. Use the derivative definition to
  establish $q'(0)$, and justify any nonexistent limit.
  \vspace*{\stretch{1}}
\end{enumerate}

\newpage
\begin{center}
  {\large\textbf{Part II --- Continued}}
\end{center}

\begin{enumerate}[leftmargin=*,itemsep=0pt,resume]
  \item Does Rolle's theorem apply to $r(x)=|x^2-1|$ on $[-2,2]$?
  Check every hypothesis. Regardless of your answer, find all
  $c\in(-2,2)$ for which $r'(c)=0$.
  \vspace{\stretch{1}}

  \item Let
  \[
    s(x)=
    \begin{cases}
      x^2, & x\le1,\\
      2x-1, & x>1.
    \end{cases}
  \]
  Verify every hypothesis of the Mean Value Theorem on $[0,3]$ and find
  all $c\in(0,3)$ satisfying its conclusion.
  \vspace{\stretch{1}}

  \item For real $a$, define
  \[
    t_a(x)=
    \begin{cases}
      x^2+2, & x\le0,\\
      ax+2, & x>0.
    \end{cases}
  \]
  Determine all values of $a$ for which the Intermediate Value Theorem,
  applied to the endpoints of $[-1,2]$, guarantees a zero in $(-1,2)$.
  For those values, find the zero and determine whether it is unique.
  Justify the theorem's hypotheses and your uniqueness claim.
  \vspace*{\stretch{1}}
\end{enumerate}

\end{document}
